\documentclass[11pt,a4paper]{article}
\usepackage{fontspec}
\setmainfont{Noto Serif}\setsansfont{Noto Sans}\setmonofont{Noto Sans Mono}[Scale=0.86]
\usepackage{amsmath,amssymb,mathtools,geometry,booktabs,tabularx,array,xcolor,hyperref,xurl,enumitem,fancyhdr}
\geometry{left=2.25cm,right=2.25cm,top=2.3cm,bottom=2.3cm}
\definecolor{ink}{HTML}{183449}\definecolor{teal}{HTML}{187C80}
\hypersetup{colorlinks=true,linkcolor=ink,urlcolor=teal,pdftitle={ECDR Week 1 Project Brief},pdfauthor={ECDR Research}}
\setlength{\parindent}{0pt}\setlength{\parskip}{5pt}\setlength{\emergencystretch}{3em}\sloppy
\setlist{nosep,leftmargin=1.5em}\newcolumntype{Y}{>{\raggedright\arraybackslash}X}
\pagestyle{fancy}\fancyhf{}\fancyhead[L]{\small\textcolor{ink}{ECDR | Week 1 Project Brief}}\fancyhead[R]{\small K=12 exploratory reference}\fancyfoot[R]{\thepage}
\begin{document}
\title{\textcolor{ink}{ECDR}\\[5pt]\Large Week 1: Project Brief and Research Grounding}
\author{ECDR Research}\date{2 November 2026}
\maketitle\tableofcontents\clearpage
% BEGIN CANONICAL WEEK1 BODY
\section{Executive brief}
This Week 1 document fixes the research question, reference contract, evidence
boundaries, and experimental plan for Event-Conditioned Certified Decision
Reuse (ECDR)
on the HAPT smartphone inertial-sensing data. The primary reference is a new,
frozen twelve-class SmallCNN1D specification. The historical six-class
checkpoint and the earlier valid-convolution/GAP theory remain provenance only.
The completed feasibility baseline is reported below; it does not establish
ECDR certification or net savings.

\subsection{Problem and decision contract}
Let $X_t\in\mathbb R^{6\times128}$ be a raw sensor window and let $P$ be
fixed per-channel z-score preprocessing, fitted on training subjects only. The
reference function is $F=f_\theta\circ P$, with logits $F(X)\in\mathbb R^{12}$
and deterministic tie-breaking. The reference decision is
$h(X)=\arg\max_k F_k(X)$. An ECDR policy may reuse a previous decision only
when its certificate proves that the current input belongs to a region on
which the same fixed reference decision holds. The primary safety target is
therefore $\hat h_{\mathrm{ECDR}}(X)=h(X)$ with zero disagreements over the
declared evaluation scope. Agreement with the reference is distinct from
agreement with ground-truth labels.

The main research question is: for the frozen SmallCNN1D K=12 reference on
HAPT, can an event-conditioned certificate preserve every reference decision
while reducing total end-to-end inference cost under a declared streaming
workload? Cost includes reference forward evaluation, preprocessing, event
detection, membership tests, certificate construction or reuse, fallback
evaluation, and I/O or scheduling overhead. Forward-only timing cannot
establish a net saving.

\subsection{Data and label contract}
The source is the UCI Smartphone-Based Recognition of Human Activities and
Postural Transitions dataset \cite{uci}, sampled at 50 Hz with three accelerometer and
three gyroscope channels. Windows are 128 samples with stride 64; each window
is tied to its recording and raw annotation coordinates. Subject-disjoint
splits are fixed as 17 training, 4 validation, and 9 test subjects. Channel
order is $[a_x,a_y,a_z,g_x,g_y,g_z]$: raw accelerometer channels are measured
in $g$ and gyroscope channels in radians per second. Windows are 2.56 seconds
long and advance by a 1.28-second step. Windows never cross recordings.
Normalizer statistics are fit across all supervised TRAIN windows and
recordings; validation and test never enter the fit. The inputs are raw sensor
signals, not the UCI 561-feature representation.

The split audit contains 9,978 full train-stream windows, 2,422 validation
windows, and 5,053 test-stream windows. The supervised
pure-or-dominant-at-least-0.6 policy contains 6,955 train, 1,751 validation,
and 3,572 test windows. The pure-basic total is 10,017 across the corpus;
these are different denominators. These counts were reproduced by the fresh
audit and recorded in the run manifest. Pure windows and
majority windows are distinct. Ambiguous or unlabeled stream windows remain
in full-stream denominators and are excluded from supervised classification
only under the declared policy. The annotation coordinate convention and the
mapping from samples to windows must be recorded with each run.

The transition snapshot has 344 pure transition windows and 10,017 pure basic
windows among 10,361 pure windows. Existing support gates therefore remain
failed for a high-quality twelve-class training claim. The explicit current
decision is to permit an exploratory K=12 feasibility reference with the
pure-or-dominant policy, while retaining the status exploratory and making no
claim that the historical G2 quality gate passed.

\subsection{Frozen reference architecture}
The K=12 contract is: input $(B,6,128)$; same-padded Conv1D
6$\to$32, kernel 5, padding 2, followed by BatchNorm, ReLU and MaxPool2;
same-padded Conv1D 32$\to$64, kernel 3, padding 1, followed by BatchNorm,
ReLU and MaxPool2; same-padded Conv1D 64$\to$64, kernel 3, padding 1,
followed by BatchNorm and ReLU; Flatten of 2,048 features; Linear
2,048$\to$128, ReLU; and Linear 128$\to$12. The executed trainable parameter
count is 283,692; shape and parameter checks passed on the instantiated K=12
model. BatchNorm evaluation uses frozen running statistics; a checkpoint
and its preprocessing parameters are an inseparable reference identity.

The operational checkpoint was selected after the completed training and
evaluation run. It uses seeds 42, 43, and 44; at most 30 epochs; Adam with learning rate
$10^{-3}$ and weight decay $10^{-4}$; cosine scheduling; and patience 10.
Epoch and seed selection use validation macro-F1. Test evaluation occurs after
selection. Required outputs include accuracy against ground truth, macro-F1,
a 12-class confusion matrix, per-class recall, reference-decision agreement,
point-margin distributions, and timing by declared scope. These outputs are
recorded for the completed feasibility run. Adjacent prediction stability is
measured for the snapshot, while ECDR/reference agreement is not established;
certificate and full ECDR cost outputs remain pending.

\subsection{Research questions and hypotheses}
\begin{description}
\item[RQ1 (reference quality)] What classification accuracy and per-class
recall does the new K=12 reference achieve against declared ground truth?
\item[RQ2 (decision preservation)] Does ECDR preserve the reference decision
with zero disagreements on stable and transition windows under the specified
certificate and lifecycle rules?
\item[RQ3 (coverage and cost)] At what event-conditioned coverage does ECDR
reduce total end-to-end cost relative to evaluating the reference on every
window, and how does that compare with a fixed fair-cost policy?
\item[RQ4 (soundness)] Which certificate bounds, membership predicates, and
lifecycle obligations are executable on this graph?
\end{description}

The hypothesis ledger is conditional. H1 predicts that stable windows will
have more stable adjacent reference predictions than transition windows; this
is not ECDR/reference agreement; the current snapshot measures it descriptively
from 2,560/2,617 stable pairs versus 23/44 transition pairs. H2 predicts that
event-conditioned certified coverage can improve at a fixed matched
total cost; this is proposed and requires a fair end-to-end benchmark. H3 is a
derived theorem obligation: if, for a candidate region $\mathcal S_r$ and
predicted class $c$, every competing margin satisfies
$\inf_{X\in\mathcal S_r}(F_c(X)-F_j(X))>0$ for $j\ne c$, then $h(X)=c$
throughout that region. This implication is mathematical, not a validated
runtime result; executable bound, arithmetic, membership, and lifecycle checks
remain pending. H4 predicts net savings only in an operating regime whose
complete cost is measured.

\subsection{Domain, event, and certificate semantics}
An anchor is a specific window $X^r$ actually evaluated by the frozen reference,
with trusted class $c=h(X^r)$. A candidate region is a set of admissible
windows, a certificate is a sound lower-bound statement over that region, and
membership is the runtime predicate that decides whether a new window belongs
to it. These are separate objects. Event descriptors are candidate signals for
choosing when to rebuild or test a certificate; they do not prove coverage by
smoothness. The whole-domain condition above is a derived condition over all
11 competing margins, not evidence that the condition holds for the current
checkpoint. Reference weights, BatchNorm running statistics, preprocessing,
tie-breaking, and arithmetic mode must remain fixed across a certificate's
lifecycle. Cached hits do not establish a new trusted anchor: reuse is allowed
only for a valid certificate and membership result, otherwise the system falls
back to the frozen reference.

\subsection{Reference classes and architecture audit}
The class IDs are fixed as follows: 1 WALKING, 2 WALKING\_UPSTAIRS, 3
WALKING\_DOWNSTAIRS, 4 SITTING, 5 STANDING, 6 LAYING, 7 STAND\_TO\_SIT, 8
SIT\_TO\_STAND, 9 SIT\_TO\_LIE, 10 LIE\_TO\_SIT, 11 STAND\_TO\_LIE, and 12
LIE\_TO\_STAND. The order is part of the checkpoint identity.

\begin{table}[ht]\centering\small
\caption{Frozen K=12 architecture audit.}
\begin{tabularx}{\linewidth}{@{}l l l r Y@{}}\toprule
Layer & Input & Output & Trainable parameters & Role\\\midrule
Input & $B\times6\times128$ & $B\times6\times128$ & 0 & Raw six-channel window\\
Conv1 (same, k5,p2) & $6\times128$ & $32\times128$ & 992 & Local temporal features\\
BN1+ReLU+Pool2 & $32\times128$ & $32\times64$ & 64 & Frozen eval affine and downsample\\
Conv2 (same, k3,p1) & $32\times64$ & $64\times64$ & 6,208 & Higher-level features\\
BN2+ReLU+Pool2 & $64\times64$ & $64\times32$ & 128 & Frozen eval affine and downsample\\
Conv3 (same, k3,p1) & $64\times32$ & $64\times32$ & 12,352 & Final temporal features\\
BN3+ReLU & $64\times32$ & $64\times32$ & 128 & Fixed evaluation normalization\\
Flatten & $64\times32$ & 2,048 & 0 & Preserve positional features\\
Linear+ReLU & 2,048 & 128 & 262,272 & Dense decision representation\\
Linear head & 128 & 12 & 1,548 & Twelve logits\\\midrule
Total & & & 283,692 & Executed trainable count; shape and parameter checks passed\\\bottomrule
\end{tabularx}\end{table}

\section{Feasibility baseline results}
The full pipeline completed with exit code 0 and 17 internal verification
checks. Validation selected seed 43 at epoch 8 by macro-F1 (0.8467853). The
selected checkpoint is \texttt{selected\_best.pt}, SHA-256
\nolinkurl{8ffc5dd8eb0e9cdf832160d074f3ae816ae6cba02adf7214a63a8b623303adb1};
the train-only preprocessing identity is
\nolinkurl{aa8e1cc20829307187e0a78200f9957f7dd068b1dc26063c2720ad832bb2fd6e}.
The three seeds selected epochs 14, 8, and 9 for seeds 42, 43, and 44;
early stopping means the histories must not be described as 30 completed
epochs.

\begin{table}[ht]\centering\small
\caption{Per-seed selection and test variability.}\begin{tabularx}{\linewidth}{@{}r Y r r r r@{}}\toprule
Seed & History length & Best epoch & Val macro-F1 & Test accuracy & Test macro-F1\\\midrule
42 & 24 & 14 & .836829 & .886618 & .767022\\
43 & 18 & 8 & .846785 & .890538 & .776265\\
44 & 19 & 9 & .836958 & .889698 & .763648\\\bottomrule
\end{tabularx}\end{table}

On the supervised test denominator of 3,572, selected seed 43 achieved
accuracy 0.8905375 and macro-F1 0.7762645. Across seeds, accuracy was
0.8889511 $\pm$ 0.0016849 and macro-F1 was 0.7689781 $\pm$ 0.0053332, where
the spreads are population standard deviations across three seeds, not
confidence intervals. The full test stream contains 5,053 windows.

\begin{table}[ht]\centering\small
\caption{Pure-plus-majority support by split; each cell is windows/subjects.}
\begin{tabularx}{\linewidth}{@{}r Y r r r@{}}\toprule
ID & Class & Train & Val & Test\\\midrule
1 & WALKING & 1077/17 & 246/4 & 531/9\\
2 & WALKING\_UPSTAIRS & 973/17 & 247/4 & 537/9\\
3 & WALKING\_DOWNSTAIRS & 894/17 & 231/4 & 492/9\\
4 & SITTING & 1103/17 & 289/4 & 546/9\\
5 & STANDING & 1206/17 & 312/4 & 597/9\\
6 & LAYING & 1188/17 & 316/4 & 586/9\\
7 & STAND\_TO\_SIT & 79/17 & 19/4 & 45/9\\
8 & SIT\_TO\_STAND & 53/17 & 13/4 & 32/9\\
9 & SIT\_TO\_LIE & 90/17 & 26/4 & 51/9\\
10 & LIE\_TO\_SIT & 89/17 & 15/3 & 46/9\\
11 & STAND\_TO\_LIE & 119/17 & 18/3 & 68/9\\
12 & LIE\_TO\_STAND & 84/17 & 19/4 & 41/9\\\bottomrule
\end{tabularx}\end{table}

\begin{table}[ht]\centering\small
\caption{Selected-seed class support, recall and F1 (supervised test, $n=3572$).}
\begin{tabularx}{\linewidth}{@{}r Y r r r@{}}\toprule
ID & Class & Support & Recall & F1\\\midrule
1 & WALKING & 531 & .9567 & .9585\\
2 & WALKING\_UPSTAIRS & 537 & .9088 & .9208\\
3 & WALKING\_DOWNSTAIRS & 492 & .9837 & .9528\\
4 & SITTING & 546 & .8407 & .8196\\
5 & STANDING & 597 & .8090 & .8429\\
6 & LAYING & 586 & .9949 & .9890\\
7 & STAND\_TO\_SIT & 45 & .7111 & .7356\\
8 & SIT\_TO\_STAND & 32 & .6563 & .7119\\
9 & SIT\_TO\_LIE & 51 & .6471 & .6168\\
10 & LIE\_TO\_SIT & 46 & .5870 & .6667\\
11 & STAND\_TO\_LIE & 68 & .5294 & .5496\\
12 & LIE\_TO\_STAND & 41 & .6585 & .5510\\\bottomrule
\end{tabularx}\end{table}

The selected-seed confusion matrix uses row and column IDs 1,2,3,4,5,6,7,8,9,
10,11,12 in the exact class order above; rows are true classes and columns are
predicted classes:
\[
\setlength{\arraycolsep}{2.5pt}
\begin{array}{r|rrrrrrrrrrrr}
&1&2&3&4&5&6&7&8&9&10&11&12\\\hline
1&508&0&22&0&0&0&0&0&0&0&0&1\\
2&16&488&15&0&0&0&2&0&0&3&2&11\\
3&4&2&484&0&0&0&2&0&0&0&0&0\\
4&0&21&0&459&57&0&1&2&1&0&5&0\\
5&0&3&0&106&483&0&1&0&1&0&2&1\\
6&0&0&0&0&0&583&0&0&0&0&3&0\\
7&1&2&0&4&1&0&32&3&1&0&1&0\\
8&0&2&0&1&7&0&1&21&0&0&0&0\\
9&0&0&0&0&1&6&0&0&33&0&11&0\\
10&0&1&0&2&0&0&0&0&0&27&1&15\\
11&0&1&1&2&0&4&2&0&20&0&36&2\\
12&0&3&2&0&0&0&1&1&0&5&2&27
\end{array}
\]

Adjacent prediction stability is descriptive only: among 5,034 valid adjacent
pairs (18 recording-boundary/gap pairs excluded), pure stable pairs agreed
2,560/2,617 (97.8219\%), pure transition pairs 23/44 (52.2727\%), and mixed
other pairs 1,399/2,373 (58.9549\%). This supports H1 descriptively for this
snapshot, but transition support is small and this is not ECDR/reference
certification or a population-level independent test.

CPU timing used one thread, FP32, batch shape $1\times6\times128$, 100 warmup
calls and 1,000 varied-input timed calls. Forward-only mean/p50/p95 were
0.173376/0.165034/0.215291 ms; reference-core preprocessing plus forward
mean/p50/p95 were 0.201103/0.183905/0.287375 ms. Both scopes exclude raw
acquisition, windowing, events, certificate construction, membership, cache
overhead and disk I/O. No speedup or net-cost claim follows. The winning point
margin median was 5.9782 and the supervised ground-truth margin median 7.5746;
these are empirical point margins, not certified bounds. The lowest selected
transition recall was STAND\_TO\_LIE, 36/68 = 0.5294.

No ECDR runtime certificate agreement has been computed. Formal certificate
agreement and full net-cost evaluation remain future work; the G2 quality gate
is HOLD because transition support is insufficient, even though the Week 1
feasibility baseline and independent baseline-artifact verification are complete.

The independent baseline artifact verifier for this Week 1 feasibility run
passed 65/65 checks with zero failures. Receipt:
\nolinkurl{output/week1\_closeout\_20261102/reference/independent\_verification.json}.
It recomputed the 131 source/data hashes, bound the 6,955 training windows and
IDs to the train-only normalizer, checked all three histories, selection,
confusion matrix, per-class precision/recall/F1 and Wilson values, reproduced
all three-seed CPU batch-512 predictions, validated CSV IDs/class order/margins,
5,034 adjacent pairs and gap mask, and checked timing metadata. The finite
batch-1 versus batch-512 check matched labels for all 5,053 stream windows with
zero disagreements and maximum absolute logit difference
$9.5367431640625\times10^{-6}$. This is an independent baseline-artifact
consistency pass, not a formal certificate verifier and not an ECDR safety
claim; the strict certificate verifier remains a Week 3 task. The G2 quality
gate remains HOLD because transition support is insufficient, and net cost
remains unestablished.

The historical parameterized \texttt{SmallCNN1D} implementation is reused to
instantiate K=12 when supplied the new class count. Its K=6 configuration and
old checkpoints remain historical and cannot transfer metrics or certificates.
BatchNorm statistics are fixed in evaluation mode; the official PyTorch 2.14
BatchNorm contract is cited in the bibliography \cite{torch_bn}.

\subsection{Objectives, split subjects, and pipeline}
The objectives are separate: (O1) measure classification accuracy and per-class
recall against ground truth; (O2) measure reference-decision agreement with a
strict zero-disagreement target; and (O3) measure net end-to-end cost against a
full-reference baseline. A single score cannot substitute for these three.

The split audit artifact \texttt{data/split\_audit.json} lists train subjects
{1,3,5,6,7,8,11,14,15,16,17,19,21,22,23,25,26}, validation subjects
{27,28,29,30}, and test subjects {2,4,9,10,12,13,18,20,24}; it reports
subject-disjoint and recording-disjoint checks with zero issues.

The reproducible pipeline is: (1) load raw HAPT recordings and annotations;
(2) construct recording-local 128-sample windows at stride 64; (3) attach raw
annotation spans and apply the declared pure/dominant label rule; (4) join the
fixed subject split; (5) fit channel means and standard deviations across all
supervised training windows only, then reuse them for validation and test;
(6) instantiate and count the K=12 graph; (7) train three seeds with the
declared protocol and select by validation macro-F1; (8) evaluate test after
selection; (9) run agreement, certificate and cost experiments separately.
The normalization fit spans training windows across records; windows themselves
never cross recording boundaries.

The sole pipeline entry point is \texttt{scripts/week1\_reference\_pipeline.py}.
The help-verified all-run command is
\texttt{./.venv/bin/python scripts/week1\_reference\_pipeline.py all --output-dir
output/week1\_closeout\_20261102/reference --epochs 30 --seeds 42,43,44 --device
cpu}. No legacy K=6 command or nonexistent K=12 config is presented as a
reproduction command. A fresh reproduction must substitute a new output
directory; the recorded directory is immutable evidence. The run record must preserve run ID, command,
environment, checkpoint hash, model config hash, preprocessing hash, data
manifest hash, and seed.

The current audit artifact is run \texttt{week1\_ref\_1793620216\_9d6c1d85},
created 2026-11-02T11:50:16Z, with pipeline hash
\nolinkurl{1ef93101f658f8a18ebf19f0c8743ec5311ffe94933aee3505be6ec29a0be4d6}.
It covers 131 source/data hashes, 122 raw sensor files, 17,453 windows from 61
recordings, float32 shape $17{,}453\times6\times128$, zero raw-slice
mismatches, zero label-builder mismatches, and verified aligned overlap for all
pairs. A one-sample offset changes 98 window policy/label assignments; this is
an explicit provenance caveat and does not change the current run. The CPU
environment is Xeon W-11855M, Python 3.11.15, PyTorch 2.14.1+cpu, NumPy 2.4.6,
pandas 3.0.6, one thread. These are audit and feasibility results, not a
quality-gate pass; ECDR certification and net-cost evaluation remain pending.

\begin{table}[ht]\centering\small
\caption{Feasibility outputs and remaining external dependencies.}\begin{tabularx}{\linewidth}{@{}lYY@{}}\toprule
Field & Evidence & Status\\\midrule
Per-class support/recall & Selected test table, all 12 classes & VERIFIED RUN\\
Confusion matrix & 12 by 12 with class IDs 1--12 & VERIFIED RUN\\
Adjacent stability & 5,034 valid pairs, stable/transition strata & VERIFIED DESCRIPTIVE\\
Reference timing & Mean/p50/p95 preprocessing plus forward & VERIFIED BASELINE\\
Three-seed variability & Selected epochs and population spread & VERIFIED RUN\\
ECDR/reference agreement & Runtime certificate comparison & NOT ESTABLISHED\\
Certificate and net cost & Formal certificate verifier and full ECDR benchmark & WEEK 3 DEPENDENCY\\\bottomrule
\end{tabularx}\end{table}

\begin{table}[ht]\centering\small
\caption{Hypothesis ledger and measurement boundary.}
\begin{tabularx}{\linewidth}{@{}lYYY@{}}\toprule
ID & Basis & Test & Status\\\midrule
H1 & Stable/transition event distinction & Adjacent reference-prediction stability; 2,560/2,617 versus 23/44, descriptive only & MEASURED DESCRIPTIVE\\
H2 & Event-conditioned reuse may avoid unnecessary rebuilds & Future matched-cost coverage, fallback rate and certificate-failure rate, stratified by pure-basic, pure-transition and mixed windows; exclude cross-recording/gap pairs & PROPOSED\\
H3 & All 11 positive regional margins imply one class & Sound bound, membership and lifecycle obligations & DERIVED\\
H4 & Reuse can save net cost in a workload regime & Full end-to-end timing including fallback and overhead & PROPOSED\\\bottomrule
\end{tabularx}\end{table}

\begin{table}[ht]\centering\small
\caption{Native Week 1 evidence register.}
\begin{tabularx}{\linewidth}{@{}lYY@{}}\toprule
ID & Scope and artifact & Status\\\midrule
W1-01 & Project brief, research grounding and targeted primary literature & VERIFIED\\
W1-02 & K=12 lock, class order and frozen checkpoint contract & VERIFIED RUN\\
W1-03 & Exact architecture graph and 283,692 trainable parameters & VERIFIED RUN\\
W1-04 & Executed full reference pipeline and reproducibility artifacts & VERIFIED RUN\\
W1-05 & Transition label policy and support blocker & VERIFIED SNAPSHOT\\
W1-06 & Main research question and separated objectives & VERIFIED\\
W1-07 & Evidence register with claim status & VERIFIED\\
W1-08 & Sensor/data full contract and denominators & VERIFIED SNAPSHOT\\
W1-09 & Hypothesis ledger H1--H4 & MEASURED/PROPOSED/DERIVED\\
W1-10 & Stable/transition decision stability plus future coverage/fallback/cost stratification & VERIFIED DESCRIPTIVE + PROPOSED\\
W1-11 & Week roadmap and dependencies & VERIFIED\\
W1-12 & English presentation and editable artifacts & VERIFIED\\\bottomrule
\end{tabularx}\end{table}

\subsection{Certificate lifecycle and architecture-specific firewall}
Reuse is permitted only when a valid certificate exists and membership holds.
Otherwise the system falls back to the fixed reference evaluation. A hit does
not renew the trusted anchor. Model or preprocessing changes invalidate the
certificate, as do recording or gap-boundary changes that violate the declared
region. The mathematical condition quantifies every point in the candidate
region and every one of the 11 competing classes; it does not assume that a
stable window is correct, and point margins are not regional bounds.

The exact overlap between adjacent stride-64 windows is 64 samples across six
channels, hence 384 shared coordinates. This is a reuse opportunity, not a
zero-difference input: the remaining coordinates change, and the full-input
distance must be computed under the declared norm if a bound requires it.
The old valid-convolution/GAP feature cancellation formulas are therefore
firewalled from this BatchNorm--MaxPool--Flatten graph. Week 3 must derive or
verify bounds on this actual graph.

\begin{table}[ht]\centering\small
\caption{Targeted related-work boundary.}\begin{tabularx}{\linewidth}{@{}lYYY@{}}\toprule
Work & Reference guaranteed & Payload/domain & Cost/evidence boundary\\\midrule
Lipschitz-margin \cite{lmt} & Perturbation margin under stated norm & Classifier robustness & No HAPT ECDR cost result\\
IBP/CROWN/auto-LiRPA \cite{ibp,crown,lirpa} & Sound bounds when relaxations and arithmetic hold & Neural-network regions & Bound runtime is method/graph dependent\\
StreamiNNC \cite{streaminnc} & Streaming CNN computation reuse & Shifted windows; padding caveats & Efficiency evidence is workload-specific\\
TinyDejaVu \cite{tiny} & Sensor-stream inference design & MCU RAM and latency & No transfer of benchmark numbers\\
LipCache/GuardNet \cite{lipcache} & Certified cache proxy boundary as defined by the work & Edge image service & GuardNet/MainNet scope must not be conflated\\\bottomrule
\end{tabularx}\end{table}

\subsection{Grounding and research gap}
The targeted primary literature covers Lipschitz-margin certification
\cite{lmt}, IBP \cite{ibp}, CROWN and auto-LiRPA \cite{crown,lirpa} for sound
bounds; StreamiNNC and TinyDejaVu \cite{streaminnc,tiny} for streaming reuse
or sensor-stream efficiency; and LipCache/GuardNet \cite{lipcache} for certified caching
boundaries. These works motivate the questions but do not establish a result
on this HAPT contract. No novelty, universal speedup, or superiority claim is
made here. A GuardNet certificate must not be described as a certificate for a
different MainNet without matching the reference boundary. The earlier
Conv-valid--ReLU--GAP derivation is architecture-specific and is not
transferred to the BatchNorm--MaxPool--Flatten reference.

\subsection{Evidence status and roadmap}
Evidence is labeled VERIFIED (read from a source artifact and checked),
DERIVED (a logical consequence whose premises remain to be checked), PROPOSED
(a planned measurement or method), or \textsc{not-established}. Week 1 closes the
brief, data and architecture scope lock, and a feasibility baseline. Week 2
documents why each data and deep-model component is used. Week 3 checks
certificate mathematics and executable bounds on the actual graph. GAP-specific
formulas from the old theory and any deep revisions are dependencies for later
weeks, not Week 1 results.

\section{Source and reproducibility contract}
Every reported number must cite a run ID, pipeline script, command, environment
and provenance hash. The closeout register records source file, scope,
denominator, status, and whether a statement is current, historical, or
proposed. The parameterized \texttt{scripts/backbone.py} implementation can
instantiate the current K=12 graph; the K=6 configuration and old checkpoints
are historical evidence only and cannot transfer metrics or certificates. User
acceptance is not asserted by this document.
% END CANONICAL WEEK1 BODY
\clearpage
\begin{thebibliography}{9}
\bibitem{uci} UCI Machine Learning Repository, \emph{Smartphone-Based Recognition of Human Activities and Postural Transitions}, 2015. \url{https://archive.ics.uci.edu/dataset/341/smartphone+based+recognition+of+human+activities+and+postural+transitions}.
\bibitem{lmt} Tsuzuku, Sato, Sugiyama, \emph{Lipschitz-Margin Training}, NeurIPS 2018. \url{https://arxiv.org/abs/1802.04034}.
\bibitem{ibp} Gowal et al., \emph{On the Effectiveness of Interval Bound Propagation}, 2018. \url{https://arxiv.org/abs/1810.12715}.
\bibitem{crown} Zhang et al., \emph{Efficient Neural Network Robustness Certification with General Activation Functions}, NeurIPS 2018. \url{https://arxiv.org/abs/1811.00866}.
\bibitem{lirpa} Xu et al., \emph{Automatic Perturbation Analysis for Scalable Certified Robustness and Beyond}, NeurIPS 2020. \url{https://arxiv.org/abs/2002.12920}.
\bibitem{streaminnc} Kechris et al., \emph{Don't Think It Twice}, 2024. \url{https://arxiv.org/abs/2408.03223}.
\bibitem{tiny} Huang and Baccelli, \emph{TinyDejaVu: Smaller RAM and Faster Inference with Neural Networks on MCUs for Sensor Data Streams}, arXiv v2, 2026. \url{https://arxiv.org/abs/2512.09786v2}.
\bibitem{lipcache} Xiang et al., \emph{LipCache: A Local Inference Proxy with Certified Caching for Edge Image Classification Service}, arXiv v1, 2026. \url{https://arxiv.org/html/2608.13144v1}.
\bibitem{torch_bn} PyTorch, \emph{BatchNorm1d}, documentation for the 2.14 API, accessed 2026-11-02. \url{https://docs.pytorch.org/docs/2.14/generated/torch.nn.BatchNorm1d.html}.
\end{thebibliography}
\end{document}

